%! Least Squares Approximations — Unit 4, Lesson 4.3 — Exit Ticket
\documentclass[10pt]{article}
\usepackage{linalg-article}
\usepackage{linalg-boxes}
\newcommand{\TallMath}[1]{$\displaystyle #1\rule[-1.4em]{0pt}{3.2em}$}
\newcommand{\xx}{\mathbf{x}}
\newcommand{\bb}{\mathbf{b}}
\newcommand{\pp}{\mathbf{p}}
\newcommand{\ee}{\mathbf{e}}
\newcommand{\T}{\mathsf{T}}

\begin{document}

\pageheader{Unit 4, Lesson 4.3}{Exit Ticket}
\namedateperiod

\vspace{0.15in}

No notes. Show your reasoning.

\begin{enumerate}[leftmargin=1.8em, itemsep=15pt]
  \item \textbf{Fit a line.} Fit $y=C+Dt$ to the points $(0,4),(1,2),(2,6)$. The set-up is
        $A=\begin{bmatrix} 1&0\\ 1&1\\ 1&2 \end{bmatrix}$, $\bb=\begin{bmatrix} 4\\2\\6 \end{bmatrix}$.
        Set up and solve the normal equations $A^{\T}A\hat{\xx}=A^{\T}\bb$, and state the best-fit line.
        \vspace{2.2cm}
  \item \textbf{What does it mean?} For your line, what is the residual $\ee=\bb-\pp$, and what does least
        squares make as small as possible? Answer in one or two sentences.
        \writelines{2}
  \item \textbf{Why not just solve $A\xx=\bb$?} Explain in one sentence why we solve the normal equations
        instead of solving $A\xx=\bb$ directly.
        \writelines{2}
\end{enumerate}

\end{document}
